\section{Conclusion and Future Work}
\label{sec:conclusion}

We presented \pimattn{}, a heterogeneous attention framework that maps decomposed asymmetric quantized attention operations to near-bank PIM for efficient LLM inference.
Our key insight is that the asymmetric dot product $Q\cdot K \approx s_Q s_K[\sum \hat{q}\hat{k} - z_K\sum\hat{q} - z_Q\sum\hat{k} + d z_Q z_K]$ decomposes into 2-bit-compatible operations ideal for near-bank PIM, while float16-intensive score@V operations remain on GPU.
Cycle-accurate simulation on Llama-3-8B demonstrates $1.4\times$ per-sequence Q@K speedup over GPU with 2-bit KV on 2048 PIM banks.
Our results show that direct 2-bit PIM compute is $2\times$ more efficient than dequantize-then-compute approaches.

Future directions include: (1) extending to prefill-mode attention where larger batch dimensions amortize asymmetric reduction overhead, (2) integrating with FlashAttention-style tiling for further memory optimization, and (3) exploring higher radix quantization (4-bit, 8-bit) with corresponding PE configurations.

\section*{Acknowledgments}
We thank the anonymous reviewers for their constructive feedback.

\newpage
\section{Appendix: Extended Results and Implementation Details}

\subsection{Complete Experiment Data}

Table~\ref{tab:full} provides the complete per-step breakdown for all six experiments.

\begin{table}[h]
\caption{Complete per-step attention latency ($\mu s$).}
\label{tab:full}
\centering
\small
\begin{tabular}{lcccccccc}
\toprule
\textbf{Config} & \textbf{Total} & \textbf{Layer} & \textbf{Gen} & \textbf{Q@K} & \textbf{Softmax} & \textbf{S@V} & \textbf{pim\_rb} & \textbf{pim\_pe} \\
\midrule
FP16 GPU   & 5623 & 176 & 95  & 97 & 0.00 & -- & -- & -- \\
FP16 PIM   & 6788 & 212 & 132 & 66 & 0.03 & 66 & 131 & 66 \\
2-bit GPU  & 2769 & 87  & 48  & 48 & 0.00 & -- & -- & -- \\
2-bit Hyb. & 3460 & 108 & 67  & 33 & 0.01 & 34 & 33  & 17 \\
2-bit All  & 3443 & 108 & 66  & 33 & 0.01 & 33 & 66  & 33 \\
2-bit Deq. & 5555 & 174 & 132 & 66 & 0.01 & 66 & 66  & 99 \\
\bottomrule
\end{tabular}
\end{table}

\subsection{Simulator Configuration}

The PIM simulator parameters used in our evaluation:
\begin{itemize}[nosep]
\item $N_{\text{banks}} = 2048$ (8 HBM3E stacks $\times$ 256 PIM banks)
\item $R = 2048$ B row buffer
\item $W_{\text{fp16}} = 16$ B/cycle, $\tau_{\text{fp16}} = 0.5$ ns (2 GHz)
\item $W_{\text{lowbit}} = 32$ B/cycle, $\tau_{\text{lowbit}} = 0.5$ ns
\item $B = 32$ GB/s per-bank DRAM-to-rowbuffer bandwidth
\item Asymmetric quant: $z_Q = 1.0$, $z_K = 2.0$
\end{itemize}

\subsection{Reproducibility}

All experiments can be reproduced using the provided simulator and configuration files.
Run \texttt{bash tools/run\_all.sh} to execute all six experiments and generate the comparison report.
The simulator source code, including the near-bank PIM latency model and per-stage timing instrumentation, is available in the repository.

